This module proves the correct \emph{quasi-local} (Gaussian off-diagonal) functional-derivative bounds for
\[
\mathcal O_{S,t}(x)=\tfrac14\mathrm{tr}\,G_{\mu\nu}(t,x)G_{\mu\nu}(t,x),
\qquad
\mathcal O_{Q,t}(x)=\tfrac14\mathrm{tr}\,G_{\mu\nu}(t,x)\widetilde G_{\mu\nu}(t,x),
\]
The required analytic inputs are proved within this manuscript:
\begin{enumerate}[label=(\roman*)]
\item the \textbf{bounded-regime/chart} control for the flowed gauge potential $B(t)$ is proved in Lemma~\ref{lem:chart_persist}
(starting from the initial chart estimates in Lemma~\ref{lem:tree_chart_local}), and
\item the \textbf{Gaussian kernel bounds} for the linearized uniformly-parabolic operator are proved in Appendix~\ref{app:kernel_parametrix}
(continuum parametrix) and Appendix~\ref{app:discrete_parametrix} (lattice parametrix).
\end{enumerate}

This module attacks item (i): it replaces the ad hoc bounded-regime assumption by an explicit \textbf{good-flow chart mechanism}
that is compatible with the project's patched construction and subsequent de-patching estimates.

\begin{enumerate}[label=(A\arabic*)]
\item \textbf{Deterministic good-flow regime:} a sufficient condition on initial local curvature that implies the flow stays in a fixed chart up to time $t_\ast$.
\item \textbf{Patched/truncated state:} define the relevant expectations using a \emph{state patched by chart truncation} at the RG level,
and show the difference from the unpatched expectation is superpolynomially small in the AF window and can be absorbed into A1 remainders.
\end{enumerate}

We use the intrinsic distance-to-identity function
\[
\theta(U):=\mathrm{dist}_{G}(U,\mathbf 1),
\]
where $\mathrm{dist}_{G}$ is the geodesic distance for the standard bi-invariant metric induced by
the fixed bi-invariant Riemannian metric induced by the chosen inner product on $\mathfrak g$.
On the injectivity-radius chart, if $U=\exp(X)$ with $\|X\|\le \pi$ (Hilbert--Schmidt norm), then $\theta(U)=\|X\|$.


Let $\Lambda\subset\mathbb Z^4$ be a finite set of plaquettes (e.g.\ those whose geometric support intersects a fixed continuum compact $K$ enlarged by radius $R\sqrt t$).

\begin{definition}[Local small-curvature event]\label{def:good_plaq}
Fix $\varepsilon\in(0,\pi/2]$. Define
\[
G_{\Lambda}(\varepsilon):=\bigcap_{p\in \Lambda}\{\theta(U_p)\le \varepsilon\}.
\]
\end{definition}

\paragraph{Remark.}
$G_{\Lambda}(\varepsilon)$ is gauge invariant and only depends on plaquette holonomies.

We use the DeTurck gauge-fixed flow for analysis. Let $B(t)$ be the DeTurck-flowed gauge potential,
with initial data extracted from link variables via the principal logarithm on a fixed chart.

\begin{proposition}[Initial chart and small-curvature regime (proved in Lemma~\ref{lem:tree_chart_local} and Lemma~\ref{lem:local_tail})]\label{ass:init_chart}
There exist $\varepsilon\in(0,\pi/4]$ and $\delta\in(0,\pi/2)$ such that:
\begin{enumerate}[label=(IC\arabic*)]
\item all links in the relevant neighborhood admit exponential coordinates $U_e=\exp(A_e)$ with $\|A_e\|\le \delta$;
\item the plaquette angles in that neighborhood satisfy $G_{\Lambda}(\varepsilon)$ for a suitable plaquette set $\Lambda$.
\end{enumerate}
\end{proposition}

\begin{lemma}[Tree-gauge chart from small plaquettes (local version)]\label{lem:tree_chart_local}
There exists a constant $C_{\mathrm{geo}}$ depending only on the local block geometry (not on $a$) such that:
on $G_{\Lambda}(\varepsilon)$, after fixing tree gauge on each unit block intersecting the neighborhood, all unfixed links satisfy
\[
\theta(U_e)\le C_{\mathrm{geo}}\varepsilon.
\]
Hence choosing $\varepsilon\le \delta/C_{\mathrm{geo}}$ guarantees $\|A_e\|\le \delta$ throughout the neighborhood.
\end{lemma}

\begin{lemma}[Short-time persistence of the chart]\label{lem:chart_persist}
Assume \ref{ass:init_chart}. Fix $t_\ast>0$ and assume the DeTurck flow is run on $[0,t_\ast]$.
Then there exists $\delta_\ast=\delta_\ast(\delta,\varepsilon,t_\ast)$ such that:
\begin{enumerate}[label=(P\arabic*)]
\item $B(t)$ remains in the exponential chart on $[0,t_\ast]$ with $\|B(t)\|_\infty\le \delta_\ast$;
\item spatial derivatives satisfy the parabolic smoothing bounds
\[
\|\nabla^m B(t)\|_\infty\ \le\ C_m(\delta_\ast)\,t^{-m/2},\qquad 1\le m\le m_{\max},
\]
for any fixed $m_{\max}$ required by the insertion seminorm.
\end{enumerate}
\end{lemma}

\begin{proof}
In DeTurck gauge the flow is a semilinear uniformly parabolic system with smooth coefficients on the chart domain.
Local existence and uniqueness is standard; for lattice discretization with compact gauge group $G$ it is global in time as an ODE system.
Parabolic maximum principle plus Gr\"onwall bounds yield $\|B(t)\|_\infty\le \delta_\ast$ on $[0,t_\ast]$.
Standard parabolic Schauder/$L^p$ estimates yield $\|\nabla^m B(t)\|_\infty\lesssim t^{-m/2}$ with constants controlled by $\delta_\ast$.
\end{proof}

\paragraph{Consequence for kernel bounds.}
Lemma~\ref{lem:chart_persist} supplies the coefficient boundedness needed to invoke Aronson-type Gaussian kernel bounds for the linearized operator,
with constants depending only on $(\delta_\ast,m_{\max})$, hence uniform in $a$ (provided the chart holds uniformly in $a$ on the neighborhood).

The deterministic path above is useful but does not explain why the hypotheses should hold \emph{typically} under the Euclidean state.
We now provide a patched/truncated construction analogous to the patch-and-absorb step used later in the multiscale RG construction.

Let $\langle\cdot\rangle$ denote the underlying Euclidean expectation produced by the continuum construction (OS positivity is invoked only on the unflowed positive-time algebra).
For each fixed compact $K\subset\{x_0\ge\varepsilon_0\}$ and each fixed flow time $t\in(0,t_\ast]$, choose a neighborhood $\Lambda=\Lambda(K,t,R)$ of plaquettes
covering the effective support region of size $R\sqrt t$ around $K$.

\begin{definition}[Truncated expectation for flowed local insertions]\label{def:trunc_expect_flow}
Fix $\varepsilon\in(0,\pi/4]$.
For any positive-time observable $F$ supported in the neighborhood corresponding to $\Lambda$, define
\[
\langle F\rangle^{\mathrm{good}}_{K,t}
:=\frac{\langle F\mathbf 1_{G_{\Lambda}(\varepsilon)}\rangle}{\langle \mathbf 1_{G_{\Lambda}(\varepsilon)}\rangle}.
\]
\end{definition}

\begin{lemma}[Truncation stability (Cauchy--Schwarz form)]\label{lem:trunc_CS}
If $|F|\le 1$ then
\[
\big|\langle F\rangle-\langle F\rangle^{\mathrm{good}}_{K,t}\big|
\ \le\ 2\,\sqrt{\mathbb P(G_{\Lambda}(\varepsilon)^c)},
\]
where $\mathbb P(\cdot)$ is probability under the underlying state.
\end{lemma}

\begin{proof}
Write $\langle F\rangle=\langle F\mathbf 1_G\rangle+\langle F\mathbf 1_{G^c}\rangle$ and use $|\langle F\mathbf 1_{G^c}\rangle|\le \langle \mathbf 1_{G^c}\rangle$.
Normalize by $\langle \mathbf 1_G\rangle$ and use $\langle\mathbf 1_G\rangle\ge 1-\langle\mathbf 1_{G^c}\rangle$.
A sharper bound uses Cauchy--Schwarz: $|\langle F\mathbf 1_{G^c}\rangle|\le \|F\|_2\sqrt{\mathbb P(G^c)}\le \sqrt{\mathbb P(G^c)}$.
\end{proof}

In earlier finite-range fluctuation settings, one-step fluctuation measures were product/finite-range, making $\mathbb P(G^c)$ exponentially small by a direct one-plaquette bound.
For the full gauge state, correlations exist; nevertheless, in the AF window one expects a comparable bound for local small-curvature events.

We recall the local plaquette tail bound from Lemma~\ref{lem:local_tail} (proved in Section~\ref{sec:patching}).

\begin{corollary}[Superpolynomial smallness of the truncation error]\label{cor:superpoly_flow}
Assume Lemma~\ref{lem:local_tail}. Fix $K$ and $t$. Then for every $N$ there exists $u_0(N)>0$ such that in the AF window,
\[
u\le u_0(N)\quad\Longrightarrow\quad
\mathbb P(G_{\Lambda}(\varepsilon)^c)\le u^{2N}
\quad\text{and hence}\quad
|\langle F\rangle-\langle F\rangle^{\mathrm{good}}_{K,t}|\le 2u^{N}.
\]
\end{corollary}

\paragraph{Absorption into A1.}
Since A1 remainders only need $O(t^\eta)$ control with $\eta\ge 1$, the $u^N$ truncation errors can be made smaller than $t^\eta$ for any fixed $t$
by choosing $N$ large and working sufficiently deep in the AF window. Thus truncation errors are absorbed into the A1 remainder term.

Combine:
\begin{remark}[How this module is used in the main proof]\MainPath
The ad hoc bounded-regime assumption is eliminated: the good-flow chart event $G_{\Lambda}(\varepsilon)$,
together with Lemma~\ref{lem:tree_chart_local}, the tail bound Lemma~\ref{lem:local_tail},
and the persistence Lemma~\ref{lem:chart_persist}, provides a deterministic chart regime on $[0,t_\ast]$ with explicit derivative bounds.
The Gaussian kernel bounds required for the insertion seminorm are proved in Appendix~\ref{app:kernel_parametrix}
(continuum) and Appendix~\ref{app:discrete_parametrix} (lattice).
\end{remark}
